\section{Discussion and Limitations}
\label{sec:reading-protocol}
The measurements establish more than the existence of a good fitted predictor, but less than a unique account of the model's computation. A computable spectral budget bounds the optimum over an entire exact orthogonal-action class. On the observed representation banks, it quantifies the additional cost of that class and permits a separate comparison with independent orthogonality. This connects a declared structural requirement to a prediction target without requiring a particular construction algorithm to attain the optimum. The constructed witness can therefore be feasible yet transfer worse than OLS without contradicting the class-risk result.

The strongest empirical lesson is that magnitude and variation matter. Positive joint lower endpoints already follow from imperfect independent regression under the theorem's assumptions. The additional group-law bounds provide more specific information, but their primary weekday sum is overwhelmingly concentrated in one condition and their upper bounds leave room for larger true costs. Likewise, the representation-type correction is essential in principle but material in only one primary condition. The complete-family deletion check preserves the primary C7 lower-bound concentration under each reduced bank's own whitening; it does not turn that concentration into evidence of uniformly large costs or replication on new contexts. Eight-coordinate results change its role and change registration rankings. These observations support a conditional characterization of the chosen spaces, not a claim that exact-action cost is uniformly large or dimension-independent.

Coherent registrations also leave plausible alternatives. A state code can realize transformations accurately and a shared invariant component can supply substantial spectral energy. Our complete small-group comparisons ask whether the prescribed image is favored within a declared finite class; even a unique spectral leader would not eliminate every semantic relabeling or identify the generating mechanism. The constructed-action and new-context comparisons add predictive evidence, but the latter concerns transport objectives rather than an independent replication of class risks or registration preferences. Its opposite domain effects further prevent treating difference fitting as a universally better objective.

The scope of the empirical conclusions follows from the observations used. Six selected models and four readouts per model provide repeated measurements, not independent samples of all language models. C7 and D4 contain small, already-seen state vocabularies; the authored-context bank extends wording, not state identities. A supplementary three-object assignment task uses the permutation group $S_3$ with new entity names and contexts; its joint held-out stratum yields only a small transport improvement over identity. Historical activation edits differ in information and intervention space, and the separate native-program branch failed clean-task competence; its confirmation remains blocked (\cref{app:behavior,app:reproduction}). Thus the present evidence does not establish successful native group computation. Finite-bank guarantees also do not supply the sampling assumptions needed for population certification.

For another representation study, the practical distinction is the object of the claim. A held-out prediction score evaluates a trained map; a class-risk bound evaluates structural compatibility under a metric; a registration margin compares coherent target actions. Keeping those objects explicit makes a negative transfer result or a competitive alternative informative rather than an apparent failure of a single recovery score.

\section{Conclusion}
We connect a type-constrained invariant-subspace budget to the minimum prediction risk of fixed-dimensional exact orthogonal actions. In frozen language-model representations, the resulting bounds quantify joint constraint costs and, with an independent orthogonal comparator, delimit the further cost of common group laws. Those guarantees are heterogeneous, constructed actions need not outperform ordinary predictors, and coherent registrations change their spectral ranking across retained spaces. New-context transport adds a separate domain-dependent boundary. Predictive recovery, exact-action compatibility and semantic preference can therefore be measured and compared without treating them as interchangeable evidence of native computation.
